\section{Limitations and future work}
\label{sec:limitations}

The fragment $\CICi$ was chosen so that every emitted axiom is provably
sound while covering the indexed phenomena the fording architecture is
about: constructor-constrained indices, forded case analysis,
index-determined singletons, recursion over families.  We delimit the
boundary, separating presentational restrictions from semantic ones,
and record the reconstruction obligations.

\subsection{Dependent index telescopes}
\label{sec:limitations-telescopes}

Index types in $\CICi$ are closed (\cref{def:env}): a family may not be
indexed over a telescope in which a later index type mentions an
earlier index ($\mathsf{dvec} : \Pi n{:}\nat.\ \Fin\,n \to \Set$).  The
restriction is inherited from the shape of source-level equations: the
forded branch equations $\tup{u} =_{\tup{\sigma}} \tup{t}_i$ are
homogeneous, and with dependent telescopes the equation at position $h$
would relate terms whose types agree only after the earlier equations
are substituted --- McBride's motivation for heterogeneous equality
\cite{McBride1999}, whose eliminator is K.  The \emph{target} side has
no such problem: untyped pointwise equations state telescopic
constraints directly, and the emitted axioms of \cref{def:theory} would
be unchanged in form.  Lifting the restriction therefore requires work
only at the source: either heterogeneous equality in the proof layer
(with its K-strength absorbed by the model, which validates it), or
cast-chained telescopic equations.  We expect no semantic obstacle ---
the stage construction and \cref{lem:index-reflection} nowhere use
closedness of index types --- but the bookkeeping in
\cref{fig:terms,fig:terms2,fig:proofs,fig:proofs2} is substantial and is left as future work.

\subsection{Polymorphism, large indices, large elimination}
\label{sec:limitations-poly}

$\CICi$ is monomorphic: type parameters are ground instances
(\cref{sec:calculus-cic}), so families indexed by \emph{types} ---
$\mathsf{eq}\ \Set\ A\ B$, universe-indexed structures --- are outside
the fragment, consistently with the intended pipeline in which
heuristic monomorphisation runs first.  Large elimination remains
excluded; fording removes its programming role in impossible branches
(\cref{ex:vhead-vtail}), the primitive $\mathsf{discr}$ covers its
role in discrimination, and its remaining uses (types computed by
recursion) would break the rigidity of set types (\cref{lem:rigid}) on
which the structural definitions of $\guardf$ and
$\semr{\cdot}{}$ rest.

\subsection{Well-founded recursion and \textsf{Acc}-style predicates}
\label{sec:limitations-wf}

The hard constraint identified in the companion paper stands, and the
indexed extension gives it a second face.  $\tfix$ captures structural
recursion only; $\mathsf{Acc}$-based recursion is excluded, because
defining equations obtained by recursion through an erased proof are
sound only with the erased propositional premises retained --- the
unconditional unfolding equation of Euclidean division at divisor $0$
is \emph{inconsistent} with $\forall n.\ n \not\feq \csS(n)$-style
facts.  On the declaration side, $\mathsf{Acc}$'s clause has its
recursive premise under a quantifier, which \cref{def:env} excludes;
semantically, its guardedness is not syntactic --- Gilbert et al.\
observe that definitional proof irrelevance for $\mathsf{Acc}$ makes
conversion undecidable \cite{GilbertEtAl2019} --- and our rank
machinery, which is the semantic image of syntactic guardedness, does
not apply.  Both faces point to the same future work: premised
unfolding equations with a semantics interpreting the defined symbol by
an arbitrary total extension.

\subsection{Acyclicity, induction, completeness}
\label{sec:limitations-completeness}

The translation emits no acyclicity axioms ($x \not\feq \csS(x)$ and
its relatives) and no induction: acyclicity is not finitely
axiomatisable over term algebras and induction is a schema
\cite{KovacsEtAl2017}; both are sound in $\ML$ (acyclicity by
\cref{cor:discrimination} and rank induction), so \emph{adding}
premise-selected instances would be sound, but the translation makes no
completeness claim that would force them.  Completeness fails further
along the axes inherited from the companion paper: the semantics is
classical and proof-irrelevant (\cref{rem:not-claimed}); conversion
under binders for lifted symbols is not reflected; goals needing
induction on families (e.g.\ $\forall n\,(v{:}\vect(n)).\
\Ihat{\mathsf{vmap}}_{\mathrm{id}}(n,v) \feq v$) need the
reconstruction backend or user induction; and expansion variants
unfold one constructor layer per occurrence, so properties requiring
unbounded inversion depth need an inductive lemma --- the same
trade-off Liquid Haskell documents for checker-driven unfolding
\cite{VazouEtAl2014}.  A completeness theorem for the ground equational
fragment over purely structural families (where \cref{prop:standard}
makes $\ML$ standard) should follow from \cref{rem:eq-shape} by
standard rewriting arguments and is left as future work.

\subsection{Functional extensionality}
\label{sec:limitations-funext}

$\ML$ is intensional (\cref{rem:funext-leak}): premise sets containing
funext-derived lemmas step outside \cref{thm:soundness}.  Repairing
this requires a coarser domain (quotient by observational equivalence),
which would disturb \cref{cor:discrimination} and the compilation
argument; nothing indexed-specific changes the analysis, but
\cref{rem:1542} raises its stakes --- the \fstar{} incident shows
extensionality axioms and index-blind encodings interacting is how
translations of this kind have actually failed.

\subsection{Reconstruction}
\label{sec:limitations-reconstruction}

Soundness here is model-theoretic; the implemented pipeline's final
arbiter is proof reconstruction, and the paper's results localise its
obligations.  Uses of inversion axioms and expanded guards replay by
\texttt{inversion}-style reasoning, derivable without axioms
\cite{CornesTerrasse1995,Monin2010} --- K-free, since the motive never
depends on the generated equations.  Uses of the unconditional
singleton-collapse equations, of deletion, and of the extra
permissiveness of \cref{def:forced-singleton} beyond CIC's parameter
mechanism (\cref{rem:singleton-criterion}) may need UIP-class tactics
at types without decidable equality; \cref{prop:guarded-scase} is the
premise-carrying alternative that eliminates this debt.  A
proof-theoretic soundness theorem in the style of \cite{Czajka2018} ---
transforming first-order derivations from $\ThE$ into $\CICi$ proof
terms modulo excluded middle, proof irrelevance and UIP --- would turn
this localisation into an algorithm and is the natural next step, with
the lifted-symbol substitution stability of the companion paper's
discussion as the known hard part.

\subsection{Other restrictions}
\label{sec:limitations-misc}

Mutual and nested inductive families, deep and mutual recursion, and
coinductive families are excluded syntactically, as in the companion
paper; each constructor of a mutual block fords independently, so we
expect the extension to be bookkeeping, but coinductive families ---
whose case-analysis content is a greatest-fixpoint inversion and whose
equality is bisimilarity --- are a genuine further step.  Enumeration-,
subset- and singleton-shaped families receive no special-cased
treatment in the formal development; the classification exists in
implementations as a choice of expansion policy, and
\cref{thm:expansion} is deliberately policy-agnostic so that any such
classification is covered.
